\documentclass[11pt,a4paper]{article}
\usepackage[T1]{fontenc}\usepackage[utf8]{inputenc}\usepackage{lmodern}
\usepackage[margin=24mm]{geometry}\usepackage{amsmath,amssymb,booktabs,longtable,microtype,xurl,hyperref}
\hypersetup{colorlinks=true,linkcolor=blue,urlcolor=blue}
\setlength{\parindent}{0pt}\setlength{\parskip}{.65em}
\title{Grothendieck Ring Classification\\Current Research Status and Audit Corrections}
\author{Fusion Frontier research project}\date{24 October 2026}
\begin{document}\maketitle
\textbf{AI-generated research; independent human verification and formalization pending.}
All theorem and classification labels describe the project's current claims.
Exact finite replays do not prove their external categorical premises, and a
verified tensor list alone does not prove enumeration completeness.

\section{Current classification scope}
Ranks two, three and four have unrestricted classification claims, with two,
five and fifteen realized based rings respectively. Higher ranks have the
explicit bounds below. Realization means existence of at least one model,
not classification of all fusion categories up to tensor equivalence.

\begin{center}\begin{tabular}{rrrrr}\toprule
Rank & Through multiplicity & Excluded & Realized & Open\\\midrule
5&20&40982&17&0\\
6&9&25362&57&2\\
7&5&9976&30&0\\
8&3&4402&93&5\\
9&2&2135&72&18\\
10&1&136&57&1\\
11&1&155&27&0\\\bottomrule
\end{tabular}\end{center}
There are \textbf{26 current generally unresolved cases}. The exact current
list and tensors are supplied in \path{data/current_remaining_cases.py}.
Rank six is complete through multiplicity three; rank eight through two;
rank nine through one. These finite statements are relative to the supplied
census and the recorded categorification decisions.

\section{R436: general inference withheld, unitary obstruction retained}
The status of \texttt{R8\_3\_436} is now unresolved for general complex
categorification, while unitary categorification remains excluded.
The proposed general obstruction needs every spherical neutral near-group
category of type $(C_3,3)$ to have a specified complete scalar associator table.
In the pinned cited source, the table's derivation assumes unitarity.
Reconstruction from abstract scalar data does not establish the required
normal form for every given spherical category. The successful exact scalar
replay cannot supply this missing categorical premise.

The revised manuscript states the scalar-table reduction conditionally and
narrows the dependent uniform near-group route claim. This correction does not
construct a category or refute the proposed nonexistence conclusion.
All other 4499 rank-eight ledger entries and all 93 models are preserved.
Exact multiplicity three has 3506 general exclusions, 16 unitary models and five
open cases. The four previously open cases retain their previously proved
sphericality and Galois-pseudo-unitarity reductions; these are not asserted for
the additional R436 case.

The correction and its precise source review are in
\path{evidence/audit_20261024/}. Earlier accepted ledgers and replay receipts
are retained as historical evidence, not silently rewritten.

\section{Other corrections}
The rank-five current replay is
\begin{verbatim}
python evidence/prepublication_20261017/rank5/replay.py --jobs 2
\end{verbatim}
It checks all 40999 inputs, 40982 general exclusions and 17 models.
The historical standalone verifier leaves the two inert-prime frontiers
unresolved in general and has complete bound seventeen. The discovery example
is a bounded attempt that may return unresolved; it is not guaranteed fresh
discovery of the stored composite certificate.

The signed rank-five commutativity lemma now explicitly requires the dimension
character to be nonzero on every basis element. The pair-product Gram matrix's
row count explicitly denotes the remaining central simple objects. The
rank-seven bibliography now names P.~G.~Budinski. The census documentation
updates the supported rank range, companion-PDF count and supplied, hash-matched
rank-seven producer source; independent full re-enumeration remains separate.
Older six-case frontier totals are scoped to their historical lower-rank range.

\section{Completed audit checks and remaining limits}
The audit checked all 1431327 tensors of its pinned census snapshot with the
supplied exact verifier and used new independent small-bound enumeration and
high-rank isomorphism programs. Its full rank-four, rank-five and rank-eleven
classification replays passed. It checked 24383 of 25362 rank-six exclusions;
979 remain outside completed replay coverage. Selected rank-seven/eight proofs
and complete hard certificates passed their stated checks. The audit's tensor
total predates the separately supplied rank-nine multiplicity-two extension.

Large highest-bound enumerations were not all independently repeated. Every
external categorical theorem and model associator was not independently
reproved. The Evans--Gannon modular-data identification used by the R820/R869
models retains an outstanding error-bound verification concern; the audit did
not establish a wrong model. No claim is made that every possible error has
been found.

\section{Reading and reproducing the atlas}
The common record directory, unresolved directory and model catalogue offer
exact structural filters for integral, weakly integral, simple, perfect,
pointed, commutative, self-dual and multiplicity-free rings. The filters concern
the underlying ring; they are independent of categorical existence. Rational
root isolation decides integrality, while complete generated-subring closure
decides simplicity. Aggregate rank--multiplicity counts retain their full
recorded populations. Definitions, examples and standard references are on the
atlas's theory page.

\begin{verbatim}
python tools/audit_corrections_20261024.py
python tools/verify_audit_revision.py
\end{verbatim}
These commands reconcile the one-record correction and validate current
publication identities and structural data. They do not claim a new full
replay of every historical categorical proof.

\url{https://sebastienpalcoux.github.io/grothendieck-ring-classification/theory.html}
\end{document}
